% ============================================================
% 正文示例章节（请将本文件内容替换为你自己的论文正文）
% ============================================================
\chapter{模板使用说明}

% 【说明】本章展示模板的常见使用方式，依据2025版编写指南编写。
% 请将本文件中所有示例内容替换为你自己的论文正文。

\section{基本排版规范}

\subsection{正文字体与段落}

正文采用小四号（12pt）字，汉字用宋体，英文和数字用 Times New Roman 体，两端对齐，段落首行左缩进2个汉字符，1.5倍行间距。以上格式由模板自动设置，写作时无需手动指定。

\subsection{标题层级}

本模板支持四级标题，格式由模板自动完成：

\begin{itemize}
    \item \textbf{章标题}：三号黑体，居中（\verb|\chapter{标题}|）
    \item \textbf{一级节标题}：四号黑体，居左（\verb|\section{标题}|）
    \item \textbf{二级节标题}：13pt 黑体，居左（\verb|\subsection{标题}|）
    \item \textbf{三级节标题}：小四号黑体，居左（\verb|\subsubsection{标题}|）
\end{itemize}

建议正文层级不超过三级节标题（即 \verb|\subsubsection|）。

\subsubsection{三级节标题示例}

三级节标题也可在必要时使用。

\section{图、表与公式}

\subsection{插图}

图题置于图片下方，使用 \verb|\caption| 命令，如\autoref{fig:sample-figure}所示。
图片文件建议放在 \texttt{figure/} 目录下。

\begin{figure}[ht]
    \centering
    \includegraphics[width=.4\linewidth]{example-image}% 替换为你的图片路径，例如 figure/your-image.png
    \caption{\label{fig:sample-figure}图片示例（请替换为实际图片）}
\end{figure}

\subsection{表格}

表题置于表格上方，使用 \verb|\caption| 命令，如\autoref{tab:sample}所示。

\begin{table}[ht]
    \caption{\label{tab:sample}自动调节列宽的表格示例}
    \begin{tabularx}{\linewidth}{|c|X<{\centering}|}
        \hline
        第一列 & 第二列 \\ \hline
        示例内容 & 示例内容 \\ \hline
        示例内容 & 示例内容 \\ \hline
    \end{tabularx}
\end{table}

\subsection{公式}

公式使用 \texttt{equation} 环境，如\autoref{equ:sample}所示。

\begin{equation}
    \label{equ:sample}
    E = mc^2
\end{equation}

\section{代码示例}

如需插入代码，使用 \texttt{lstlisting} 环境，如下所示。

\begin{lstlisting}[%
    language={C},
    caption={示例代码},
    label={code:sample}
]
#include <stdio.h>

int main(void)
{
    printf("Hello, Westlake University!\n");
    return 0;
}
\end{lstlisting}

\section{引用}

\subsection{文献引用}

使用 \verb|\cite{key}| 引用参考文献，例如：某某研究\cite{sample-article}表明……

参考文献条目在 \texttt{body/ref.bib} 文件中以 BibTeX 格式维护。

\subsection{图、表、公式的交叉引用}

使用 \verb|\autoref{label}| 进行交叉引用，模板已预设中文名称（图、表、式等）。

\section{脚注}

脚注内容采用小五号字\footnote{这是一个脚注示例。中文用宋体，英文和数字用 Times New Roman 体。}。

\chapternonum{同一页上的章标题}

\verb|\chapternonum{...}| 用于生成不带编号的章级标题（如参考文献、致谢等），该命令同时将其加入目录并设置篇眉。
